\documentclass[11pt,a4paper]{article}
\usepackage[T1]{fontenc}
\usepackage[utf8]{inputenc}
\usepackage{amsmath,amssymb,amsthm}
\usepackage[margin=1in]{geometry}
\usepackage[colorlinks=true,linkcolor=blue,urlcolor=blue]{hyperref}
\theoremstyle{plain}
\newtheorem{theorem}{Theorem}
\newtheorem{lemma}[theorem]{Lemma}
\newtheorem{proposition}[theorem]{Proposition}
\newtheorem{corollary}[theorem]{Corollary}
\theoremstyle{definition}
\newtheorem{remark}[theorem]{Remark}

\title{The empirical recurrence for OEIS A197230 is false, and the correct one, of order 25}
\author{Adrian Perez Fontelles\\ \small Independent researcher}
\date{1 September 2026}
\begin{document}
\maketitle

\begin{abstract}
OEIS A197230 carries an empirical recurrence of order $22$ contributed by R.~H.~Hardin. It is
false. The entry publishes $22$ terms and has offset $1$, so the recurrence first says
something at $n=23$, and there it already fails: the true value is
$1327965802062332$ and the recurrence predicts $1327965802062198$, a difference of $134$. The
sequence is a transfer-matrix count, so it is $C$-finite and its true minimal recurrence can
be computed and proved; it has order $25$, and its first $22$ coefficients are exactly the
published ones. The published line is therefore the correct recurrence with its last three
terms dropped. Both statements --- the failure and the corrected recurrence --- are proved
here in exact integer arithmetic.
\end{abstract}

\noindent\small 2020 Mathematics Subject Classification. 05A15, 11B37, 15B36.\normalsize

\section{The sequence and the claim}

OEIS A197230 is ``Number of n X 3 0..4 arrays with each element x equal to the number its horizontal and vertical neighbors equal to 4,3,0,1,2 for x=0,1,2,3,4.''. It has offset $1$ and its published DATA is
\[
2, 12, 55, 242, 1178, 5355,\ \dots,\ 283580433219551
\]
--- $22$ terms in all. The entry carries
\begin{quote}\small
Empirical: a(n) = 2*a(n-1) +3*a(n-2) +38*a(n-3) +40*a(n-4) +10*a(n-5) -150*a(n-6) -375*a(n-7) -175*a(n-8) +637*a(n-9) +2013*a(n-10) +1368*a(n-11) -1614*a(n-12) -4361*a(n-13) -4043*a(n-14) +946*a(n-15) +4156*a(n-16) +3551*a(n-17) +552*a(n-18) -881*a(n-19) -629*a(n-20) -1033*a(n-21) +508*a(n-22).
\end{quote}
As of the ``Last modified'' line on the live entry (Feb 07 2026 20:53:17, revision
11) the statement is recorded as empirical and nothing on the entry records it
as settled or as corrected. Written out, the claim is
\begin{equation}\label{eq:conj}
a(n)\;=\;2\,a(n-1) + 3\,a(n-2) + 38\,a(n-3) + 40\,a(n-4) + 10\,a(n-5) - 150\,a(n-6) - 375\,a(n-7) - 175\,a(n-8) + 637\,a(n-9) + 2013\,a(n-10) + 1368\,a(n-11) - 1614\,a(n-12) - 4361\,a(n-13) - 4043\,a(n-14) + 946\,a(n-15) + 4156\,a(n-16) + 3551\,a(n-17) + 552\,a(n-18) - 881\,a(n-19) - 629\,a(n-20) - 1033\,a(n-21) + 508\,a(n-22),
\end{equation}
of order $22$, with no range stated.

Note at once that \eqref{eq:conj} is untestable on the published terms: with offset $1$ and
order $22$ the first index at which it asserts anything is $n=23$, and the entry stops at
$n=22$.

\section{The model}

The condition is local to a cell and its four horizontal and vertical neighbours. Write
$x_{t,u}$ for the entry in row $t$ and column $u$, $1\le u\le3$, with values in
$\{0,1,2,3,4\}$, and let
\[
f=(4,3,0,1,2),\qquad\text{that is } f(0)=4,\ f(1)=3,\ f(2)=0,\ f(3)=1,\ f(4)=2 .
\]
The entry asks that every cell satisfy
\begin{equation}\label{eq:loc}
\#\bigl\{(d_1,d_2)\in\{(0,-1),(0,1),(-1,0),(1,0)\}:\ x_{t+d_1,\,u+d_2}=f(x_{t,u})\bigr\}
\;=\;x_{t,u},
\end{equation}
positions outside the array not being counted. Because the neighbourhood reaches one row up
and one row down, a state carrying a single row is not enough: take ordered pairs of
consecutive rows as vertices, $|\mathcal V|=S=15625$, with $(p,c)\to(c,x)$ an edge when every
cell of the middle row $c$ satisfies \eqref{eq:loc} with $p$ above and $x$ below, and mark
the pairs whose first row passes with nothing above, and those whose second row passes with
nothing below. Then $a(n)=\mathbf{v}^{\!\top}N^{\,n-2}\mathbf{u}$ for $n\ge2$.

This model reproduces the entry's DATA at all $22$ published indices exactly. It was also
checked against a direct enumeration of all $5^{3n}$ arrays for $n=1,2$, which gives $2$ and
$12$, the entry's first two terms.

\section{The published recurrence fails}

Evaluating the model at $n=23$ gives
\[
a(23)\;=\;1327965802062332,
\]
while the right-hand side of \eqref{eq:conj}, applied to the $22$ published terms, gives
\[
1327965802062198.
\]
The two differ by $134$. Since \eqref{eq:conj} is asserted without any range and
$n=23$ is the first index at which it asserts anything, the published recurrence is false.
It continues to fail at $n=24,25,26,\dots$; nothing about $n=23$ is exceptional.

\section{The recurrence that does hold}

Being a walk count, $a$ is $C$-finite. Solving for a constant-coefficient recurrence over
$\mathbb{Q}$ from the model's own terms, the least order that admits one is $25$, and the
coefficients come out integral:
\begin{equation}\label{eq:true}
a(n)\;=\;2\,a(n-1) + 3\,a(n-2) + 38\,a(n-3) + 40\,a(n-4) + 10\,a(n-5) - 150\,a(n-6) - 375\,a(n-7) - 175\,a(n-8) + 637\,a(n-9) + 2013\,a(n-10) + 1368\,a(n-11) - 1614\,a(n-12) - 4361\,a(n-13) - 4043\,a(n-14) + 946\,a(n-15) + 4156\,a(n-16) + 3551\,a(n-17) + 552\,a(n-18) - 881\,a(n-19) - 629\,a(n-20) - 1033\,a(n-21) + 508\,a(n-22) + 28\,a(n-23) - 82\,a(n-24) + 40\,a(n-25).
\end{equation}
The first $22$ coefficients of \eqref{eq:true} are, term for term, those of
\eqref{eq:conj}. The published line is \eqref{eq:true} with its last three terms,
$+28\,a(n-23)-82\,a(n-24)+40\,a(n-25)$, missing.

\begin{theorem}
\eqref{eq:true} holds for every $n>26$.
\end{theorem}

\begin{proof}
Let $q(t)=t^{25}-\sum_i c_it^{25-i}$ be its characteristic polynomial and
$w=q(N)\mathbf{u}$. For $n\ge27$,
\[
a(n)-\sum_i c_i\,a(n-i)=\mathbf{v}^{\!\top}N^{\,n-27}w ,
\]
so the recurrence holds from there on if and only if $\mathbf{v}^{\!\top}N^{m}w=0$ for
every $m\ge0$; by the Cayley--Hamilton theorem $N^{S}$ is an integer combination of
$I,N,\dots,N^{S-1}$, so checking $m<S$ suffices. The vector $w$ was formed by $25$
matrix--vector products in exact integer arithmetic and the run continued until $S+1$
consecutive zeros had been seen. Every value vanished from the index corresponding to $n=27$
onwards.
\end{proof}

Solving for the recurrence is only a way of PROPOSING \eqref{eq:true}; what makes it a
theorem is the annihilation test above, which is independent of how the candidate was found.

\section{Verification}

Four checks.

First, the model reproduces all $22$ published terms exactly, and its first two values were
confirmed by brute-force enumeration over all $5^{3n}$ arrays.

Second, the failure at $n=23$ was computed twice: once by iterating the transfer matrix, and
once by applying \eqref{eq:conj} to the entry's own published integers with no matrices
involved. Both are exact integer computations and they disagree by $134$.

Third, \eqref{eq:true} was checked directly on the model's terms for $n=27,\dots,120$
before any matrix argument was made, and holds at every one.

Fourth, the annihilation test was carried to the full Cayley--Hamilton bound rather than to a
sampled prefix.

\begin{remark}
The shape of the error is worth stating plainly, because it decides what this paper is. The
published coefficients are not wrong; they are incomplete. A recurrence fitted to a finite
sample and then reported at an order the sample cannot support will look like this: correct
as far as it goes, and false at the first index where it says anything. The correction is
\eqref{eq:true}.
\end{remark}

\begin{thebibliography}{9}
\bibitem{oeis} The OEIS Foundation, \emph{The On-Line Encyclopedia of Integer
Sequences}, \url{https://oeis.org}, 2026. Sequence A197230.
\bibitem{stanley} R.~P.~Stanley, \emph{Enumerative Combinatorics}, Volume 1, 2nd ed.,
Cambridge University Press, 2011. (Transfer-matrix method, Section 4.7.)
\bibitem{kp} M.~Kauers and P.~Paule, \emph{The Concrete Tetrahedron}, Springer, 2011.
($C$-finite sequences and their recurrences.)
\bibitem{horn} R.~A.~Horn and C.~R.~Johnson, \emph{Matrix Analysis}, 2nd ed., Cambridge
University Press, 2013.
\end{thebibliography}

\end{document}
